\documentclass{article}

\usepackage[a4paper]{geometry}
\usepackage[english]{babel}
\usepackage{parskip}

\usepackage{amsmath, amssymb}
\usepackage{graphicx}
\usepackage{subcaption}
\usepackage{float}
\usepackage[hidelinks]{hyperref}

\usepackage{listings}
\usepackage{xcolor}

\definecolor{codegreen}{rgb}{0,0.6,0}
\definecolor{codegray}{rgb}{0.5,0.5,0.5}
\definecolor{codepurple}{rgb}{0.58,0,0.82}
\definecolor{backcolour}{rgb}{0.95,0.95,0.92}

\lstdefinestyle{mystyle}{
  backgroundcolor=\color{backcolour},
  commentstyle=\color{codegreen},
  keywordstyle=\color{magenta},
  numberstyle=\tiny\color{codegray},
  stringstyle=\color{codepurple},
  basicstyle=\ttfamily\footnotesize,
  breaklines=true,
  captionpos=b,
  keepspaces=true,
  numbers=left,
  numbersep=5pt,
  showstringspaces=false,
  tabsize=2
}
\lstset{style=mystyle}

\newcommand{\PlotDir}{Plots}

\title{NURb Hand-in 1}
\author{Milan Tol (3162168)}
\date{Date}

\begin{document}
\maketitle

\textit{Write a short description of your solution in each section. Explain your choices. Draw the conclusions from your results.}

\textit{You can refer to selected lines of code where you think it is necessary. Make sure the line numbers you reference are after ``black'' has been run on your files.}

\textit{Don't forget to describe your figures and mention what conclusions you draw from them!}


\section{Simulating the solar system}


\subsection{(1a)}

In this exercise we will simulate the trajectory of objects in the solar system. 
We obtain the initial conditions from astropy. These include the positions in Mpc and velocities in Mpc/yr
for the sun and all 8 planets. The positions in the xy and xz plane can be seen in Figure \ref{Fig:initial_conditions}.
\begin{figure}[H]
    \centering
    \includegraphics[width=0.85\textwidth]{\PlotDir/initial_positions.png}
    \caption{Initial positions of the objects in the solar system in the x-y and x-z planes.}
    \label{Fig:initial_conditions}
\end{figure}


\subsection{(1b)}

Show the orbits of all objects in the x - y plane
in a suitable reference frame, as well as a plot of time versus position in the z plane. 
Explain whether leapfrog is a suitable choice of an algorithm, and why.

Next, we compute the orbits over the following 300 years from these initial conditions. 
We do this by solving the differential equation given by Newton's law of gravitation. 
Since this system is isolated, it conserves energy. Because the system satisfies this property,
it is best to use an integrator that satisfies time-symmetry since this implies the 
conservation of energy. The \textit{leapfrog integrator} is such a ODE-solver and hence is a 
suitable choice for the problem at hand. 
It works by alternatively updating the position and velocity, rather than updating them simultaneously. 
This means that it uses the velocity at time step $t+dt/2$ to compute the change in position from 
$x(t)$ to $x(t+dt)$. 
It then proceeds to use the position at $x(t+dt)$ to update the velocity from 
$v(t+dt/2)$ to $v(t+3dt/2)$. However, our initial conditions return the velocity and position at some 
starting time $t_0$. So to start this process we must first update the velocity by half a time step
by using a different method. We perform this initial \textit{kick} using the 4th order Runge-Kutta method.
This method does not respect energy conservation, but since we do it for a half time-step, 
we lose a minimal amount of accuracy. 

For this computation we use a time step of 0.8 days.
The orbits of each object in the xy plane using leapfrog can be seen in Figure
\ref{fig:orbits_xy_leapfrog}. One can see ellipses in the Figure. This is because after some 
period T, each object returns approximately to its original state in phase-space. 
In a general multi-body problem this is unexpected, but for the 2-body problem this is a 
well-known result, so why does it apply here? This is because the sun dominates the mass-budget
of the system. Thus the system is well approximated by several 2-body problems of planet-sun pairs.
The z-coordinate (the elevation from the plane) versus time can be seen in Figure
\ref{fig:orbits_xy_leapfrog}. We see oscillatory motion for all objects. 
This is due to the same reason as mentioned above for the xy plane. From this 
plot we can read off the period $T$: the time it takes for a body to do one sinusoidal 
oscillation in the z-axis.

\begin{figure}[H]
    \centering
    \includegraphics[width=0.85\textwidth]{\PlotDir/orbits_xy_leapfrog.png}
    \caption{Orbits of the objects in the x-y plane using the leapfrog integrator over a period of 300 years.
    The timestep for the integrator is 0.8 days. Note that all
    objects approximately return to their starting positions after some period, resulting in ellipses.
    This is expected in a system where one mass, the sun, dominates.}
    \label{fig:orbits_xy_leapfrog}
\end{figure}
 
\begin{figure}[H]
    \centering
    \includegraphics[width=0.85\textwidth]{\PlotDir/z_vs_time_leapfrog.png}
    \caption{Position in the z plane as a function of time for all objects using the leapfrog integrator over a period of 300 years.
    The timestep for the integrator is 0.8 days. Note that all
    objects approximately return to their starting positions after some period, resulting in oscillatory motion.
    This is expected in a system where one mass, the sun, dominates.}
    \label{fig:orbits_zt_leapfrog}
\end{figure}

    
\subsection{(1c) }

In this subsection, we compute the orbits again over a period of 300 years. But this time, we choose 
a less suitable algorithm that does not respect energy conservation.
We will use 4th order Runge-Kutta (RK4) throughtout the entire period. 
Since this method is a forward integration method, it will likely artificially induce
energy into the system due to finite stepsizes. Thus we expect the orbits to grow larger over time. However, the 
further the objects are away from the sun, the typical time scales corresponding to their orbit, their period $T$, 
will get increasingly large relative to the timesteps used for integrating. Effectively, 
our time steps become smaller compared to the timescale of the dynamics of the object.
Hence RK4 becomes more accurate the further separated the body is from the sun. 

We test our hypothesis by analysing Figure \ref{fig:orbits_xy_RK4} and \ref{fig:z_vs_time_RK4}.
We see that in the xy plane, the gas giants that are located at larger radii traverse ellipses 
and seem to approximately return to their initial phase-space states. For the inner 4 planets, 
we see a clear difference between the trajctory computed by Leapfrog and RK4.
To better see what is happening to these trajectories we plot a zoomed-in version of the 
trajectories of the sun and the inner 4 planets in Figure \ref{fig:orbits_xy_RK4_inner}. 
We see that the planets do not return to their initial states. Instead, the planets migrate further out
as we would expect if the integrator artificially added energy into the system. We also see that this
effect is stronger for the bodies closer in, mercury in particular.

We can confirm that they are indeed migrating outwards by looking at Figure \ref{fig:z_vs_time_RK4} 
and Figure \ref{fig:r_difference_RK4_minus_leapfrog}. In \ref{fig:z_vs_time_RK4} we see that the 
amplitude of the z_separation oscillation of the terrestrial planets increases over time. Hinting 
at outwards motion. In \ref{fig:r_difference_RK4_minus_leapfrog} we plot the difference in radial distance of 
bodies of the trajectories computed by RK4 versus leapfrog. We normalize it by the max radius of the trajectories 
found by using leapfrog. This gives us a better comparison of objects close and further from the sun, to see where the RK4 
algorithm fails more. 
We see that for the terrestrial planets there is a clear trend that the bodies are moving outwards. 
For the gas planets, there is a suprising amount of deviation as well, but it seems to oscillate around 0.
Hence there is no substantial amount of energy added for these bodies further out.
We have checked that decreasing the step-size increases the deviation strongly and  
jupiter starts showing outwards migration as well.

\begin{figure}[H]
\centering
\includegraphics[width=0.8\textwidth]{\PlotDir/orbits_xy_RK4.png}
\caption{Orbits of the objects in the x-y plane using a 4th order Runge-Kutta integrator over a period of 300 years.
    The timestep for the integrator is 0.8 days.}
\caption{fig:orbits_xy_RK4}
\end{figure}

\begin{figure}[H]
\centering
\includegraphics[width=0.8\textwidth]{\PlotDir/orbits_xy_RK4_inner.png}
\caption{Orbits of the inner 4 planets and the sun in the x-y plane using a 4th order Runge-Kutta integrator over a period of 300 years.
    The timestep for the integrator is 0.8 days.}
\caption{fig:orbits_xy_RK4_inner}
\end{figure}

\begin{figure}[H]
    \centering
    \includegraphics[width=0.8\textwidth]{\PlotDir/z_vs_time_RK4.png}
    \caption{Position in the z plane as a function of time for all objects using a 4th order Runge-Kutta integrator 
    over a period of 300 years. The timestep for the integrator is 0.8 days.}
    \label{fig:z_vs_time_RK4}
\end{figure}

\begin{figure}[H]
    \centering
    \includegraphics[width=0.8\textwidth]{\PlotDir/r_difference_RK4_minus_leapfrog.png}
    \caption{Difference in the x positions of the bodies as a function of time between the leapfrog method and another integration method.}
    \label{fig:r_difference_RK4_minus_leapfrog}
\end{figure}

\subsection{(1d)}

The movie has a framerate of 24 per second. Note that the period of venus is about the same 
as the amount of time that passes in between frames. This results in a seemingly 
long period, but is not accurately depicting the actual period. 


\section{Calculating potentials in different ways}

\subsection{(2a) }

Describe how you constructed your octree.

The octree was constructed by recursively splitting nodes into their children and passing the information to them.
For this we first constructed a Node class. This node class stores the following information:
\begin{itemize}
    \item Indices of the particles inside the node.
    \item The depth/level of the node. The root node has depth/level 0. 
    \item The index of the node. The index is defined as the position of the node, in units of nodes at the level of that node.
    In the first split: the root node splits into 8 children of depth 1. The child that is closest to the origin in position space
    gets index (0,0,0), while the node the furthest away gets index (1,1,1). In the next split (depth 2) the furthest node would instead 
    have index (3, 3, 3). 
    \item The center of mass of the node. 
    \item The children of the node, which is an array containing 8 nodes. If the node is a leaf node, then it contains no children.
    \item The total mass of the Node. Since we use a constant particle mass this is simply the number of particles times the particle mass.
\end{itemize}

The octree was the constructed by recursively calling the \textit{build} function. 
The build function is passed an array of particle indices, particle positions, the center of the current node,
the index of the current node, the box_size of the current node, depth of the current node, and the desired max depth of the tree.
The build function essentially does the following: It checks whether the current Node 
contains any particles or whether the max desired depth of the tree is reached. If that is the case, then it 
returns the current Node without splitting it up into any children. If neither condition is met, the build function 
splits up the current node into 8 children and then calls the build function on each child before returning the 
current Node. Note that it is important that we pass the box_size of the current node to calculate the center_positions 
of the child nodes. 
So we can construct our tree by creating a root node with the build function on 
all our particles and pass a desired depth.

The mass map at some desired resolution can then be constructed by counting 
the particles in each node at the level corresponding to the desired resolution.
More technically: For each pixel in the mass map, we call the get_node_at_level function.
This function traverses the octree and returns the node at a given level and index. 
The massmap value is then given by the total mass in the node.
Traversing the octree can be done efficiently since the index of the node has a neat mapping 
to all the splits of parent nodes necessary to get to the desired child. 
See the comments provided in the get_node_at_level function.

The mass distribution for four slices at different depth levels, namely the first four slices in the x-axis, at levels 3,
5 and 7, are plotted in figures \ref{Fig: massmap level 3, Fig: massmap level 5, Fig: massmap level 7}.

\begin{figure}[H]
    \centering
    \includegraphics[width=0.8\textwidth]{\PlotDir/fig2a_level3.png}
    \caption{Mass distribution for the first four slices in xi at level 3 of the octree.}
    \label{Fig: massmap level 3}
\end{figure}

\begin{figure}[H]
    \centering
    \includegraphics[width=0.8\textwidth]{\PlotDir/fig2a_level5.png}
    \caption{Mass distribution for the first four slices in xi at level 5 of the octree.}
    \label{Fig: massmap level 5}
\end{figure}
    
\begin{figure}[H]
    \centering
    \includegraphics[width=0.8\textwidth]{\PlotDir/fig2a_level7.png}
    \caption{Mass distribution for the first four slices in xi at level 7 of the octree.}
    \label{Fig: massmap level 7}
\end{figure}

As expected we see that the lower levels return a lower resolution. The spatial fraction that each node
occupies is equal, this results in an uneven spread of particles that each node contains, as expected 
for an octree.
        

\subsection{(2b) }

Describe how you computed the potential and implemented fft.

To compute the potential we can transform the differential equation into fourier space. 
But this requires us to know the fourier transform of the density. 
For this we first grid the particles into a $128^3$ grid. This grid is given by level 7 nodes of our octree.
Again we can create a massmap that we can convert into a density grid by dividing it by the boxsize of level 7 nodes.
The volume of a level 7 node is given by $\left( \frac{L}{128}\right)^3$.
This 3 dimensional grid containing the densities is then passed into our 3d fast fourier transform (fft) algorithm.
The 3d fft is essentially just a bunch of 1d ffts along each axis combined. The interesting 
part is programming the 1d fft. We do this using the Cooley-Tukey algorithm with a bottom up approach.

We start from the fourier transforms of single array elements, which are simply those array elements.
Then we combine them by treating them as the even and odd terms of an array containing two elements
as prescribed by the Danielson-Lanczos lemma. 
We keep adding these terms with the appropriate exponentials until we end up with 
the fourier transform of the full array. 
In order to save us from computing all the exponentials, we use trigonometric recurrence.
This allows us to compute the exponential values by applying products and summing, which 
are operations that require significantly less flops than complex exponentiation. 

After having performed the ffts in each axis, we end up with the 3d fourier transform. 
The fourier transform of the potential is then computed by dividing the modulus squared of 
the wavevector $k$, and multiplying with $-\frac{G}{\pi}$. 
To obtain the potential in real space, we must take the inverse fourier transform. 
In our definition this only requires us to add a minus sign in the complex exponent of our 
fft algorithm. Ultimately the result is independent of the definition of the fourier transform used.
We plot the resulting potential in Figure \ref{fig:fig2b} in various slices of the x-axis. 
Note that we plot the negative of the potential, so yellow spots correspond to deeper wells which 
we expect to be located in denser environments, while blue/black spots correspond to underdense voids.
The way the discrete fourier transform is formulated,
assumes a periodicity in the axis over which the transform is done. 
The periodicity of the fourier transform is apparent in the plot. One can see that 
in particular the voids connect from one edge to the other.  

When comparing the massmap in Figure \ref{Fig: massmap level 7} with the potential in Figure \ref{fig:fig2b}
in the 0th x slice,
we see that the deep well at y = 110 and z = 10 corresponds to a lot of mass that is accumulated.
It is hoever suprising that we see a void in the potential at around y=50 and z=80. In the massmap 
we see that there is quite some structure there. The discrepancy can be explained by the fact that our plots 
are 2d. The potential depends however on structure in different x-slices as well. At different x-slices there 
might be a lot less mass in the y,z position.


\begin{figure}[H]
    \centering
    \includegraphics[width=0.8\textwidth]{\PlotDir/fig2b.png}
    \caption{Negative of the potential computed using the FFT method for different slices in the x-axis.
    Yellow spots correspond to deeper wells which 
    we expect to be located in denser environments, while blue/black spots correspond to underdense voids.}
    \label{fig:fig2b}
\end{figure}


\section{Spiral and elliptical galaxies}

\subsection{(3a) }

Plot the distributions of the rescaled features, and save them to a text file; print the features
for the first 10 objects in your report (do not hardcode them into the report, but show the first
10 lines of the text file).

For the data processing we have taken additional steps besides feature scaling to ensure 
the machine can learn more optimally. When inspecting the data with 
simple feature scaling (see Figure \ref{fig3a_old}) we see that 
the extendedness as well as the emission line flux appear to be log distributed. 
This can cause problems since the features of the bulk of objects around the mean 
are then indistinguisable for the logistic regression algorithm. 
To ensure that our algorithm can still extract information from these 2 features, we 
transform them into log space before applying the feature scaling.
We also see that the emission line flux suffers from 2 extreme outliers. 
Hence without additional preprocessing, the line emission of other objects is still indistinguishable.
So we also clip the emission line flux after transforming the data into logspace as follows:
We demand that each data set is within some median offset from the median line emission flux.
In this case we set this tolerance to 100 median offsets. If the data falls outside of this range, 
it is clipped to 100 median offsets from the median. The distribution of the preprocessed 
data can be seen in Figure \ref{fig3a}. We see that 5 line emission data points are clipped 
and there is now a visible distribution in the bulk of line emission data points. 
Also the values of the extendedness of the galaxies are more distinguishable.

\begin{figure}[H]
    \centering
    \includegraphics[width=0.8\textwidth]{\PlotDir/fig3a_old.png}
    \caption{Distributions of the rescaled features for the spiral and elliptical galaxies without any 
    additional preprocessing. Note that the y-axis for the extendedness and the emission line flux are logscaled.}
    \label{fig3a_old}
\end{figure}
\begin{figure}[H]
    \centering
    \includegraphics[width=0.8\textwidth]{\PlotDir/fig3a.png}
    \caption{Distributions of the rescaled features for the spiral and elliptical galaxies.  Note that the y-axis for the extendedness and the emission line flux are logscaled.
    Also the distribution in the line emission is clipped and the plotted values for the extendedness 
    and emission line flux are logarithmic.}
    \label{fig3a}
\end{figure}


    \lstinputlisting[
        firstline=1,
        lastline=11,
        basicstyle=\ttfamily\small,
        frame=single
    ]{\PlotDir/galaxy_data_scaled.txt}

\subsection{(3b) }

Explain how you implemented logistic regression. 
 Plot the evolution of your cost function with the number of iterations for different pairs of features.
 Explain how does it change depending on the feature pair.

\begin{figure}[H]
    \centering
    \includegraphics[width=0.8\textwidth]{\PlotDir/fig3b.png}
    \caption{Evolution of the cost function with the number of iterations for different pairs of features in the logistic regression model.}
    \end{figure}


\subsection{(3c) }

Print the values of true/false positives/negatives, as well as the F1 score.
\lstinputlisting[language=Python]{\PlotDir/logistic_regression_metrics.txt}

Plot the two features against each other and indicate the decision boundary and discuss your results.
\begin{figure}[H]
    \centering
    \includegraphics[width=0.8\textwidth]{\PlotDir/fig3c.png}
    \caption{Plot of the pairs of two features against each other with the decision boundary.}
    \end{figure}
\end{document}